\section{Multi-Person Human Pose Estimation}
\label{sec_lit_hpe}

Human pose estimation (HPE) is the problem of localising a fixed set of
anatomical keypoints - such as the seventeen joints of the COCO
skeleton~\cite{lin2014coco} - on every person present in a given image.
Given the breadth of the domain, this work is restricted 
to 2D, multi-person estimation.
In this context the goal is to find per-image keypoints
for an unknown number of people without the use of temporal context
or 3D reconstruction.
The following three useful surveys provide considerable background and scope:
Chen et al.~\cite{chen2022hpesurvey} cover monocular HPE end-to-end,
Zheng et al.~\cite{zheng2023hpesurvey} provide a more recent
deep-learning-focused survey, and Andriluka et al.~\cite{andriluka2014mpii}
anchor the modern benchmark culture with the MPII dataset.

The single-person pose estimators that underpin most multi-person systems
emerged from a series of convolutional architectures.
Toshev and Szegedy's DeepPose~\cite{toshev2014deeppose} (a seminal work) framed
keypoint localisation as direct coordinate regression from a deep
network.
Subsequent architectures shifted to dense heatmap prediction:
convolutional pose machines~\cite{wei2016cpm} introduced iterative
belief-map refinement, and the stacked hourglass
network~\cite{newell2016hourglass} established the
encoder-decoder-with-skip-connections topology that underpins most
modern single-person estimators.
Xiao et al.'s ``simple baselines'' work~\cite{xiao2018simplebaselines}
showed that a ResNet trunk with three deconvolution layers matched
hourglass accuracy at lower complexity (an important finding given the trend of increasing model size).
The HRNet family~\cite{sun2019hrnet} pushed this further by running
several parallel branches at different spatial resolutions throughout
the network and exchanging information between them at every stage,
rather than recovering resolution through up-sampling at the end of an
encoder-decoder pass.
The resulting access to high-resolution feature maps yielded
sharper keypoint heatmaps and lower localisation error than any prior
backbone of comparable parameter count, which is why HRNet became the
gold-standard backbone for most top-down systems and for the bottom-up
HigherHRNet variant~\cite{cheng2020higherhrnet} discussed in
Section~\ref{sec_lit_grouping}.

\subsection{Top-down approaches}
\label{sec_lit_hpe_topdown}

Top-down systems handle multi-person HPE by
first detecting each person (using a person detector)
and then applying a single-person pose estimator to each detected bounding box.
This architecture is most notably represented by Mask R-CNN~\cite{he2017maskrcnn},
which extends the Faster R-CNN detector with a per-keypoint mask head,
and AlphaPose~\cite{fang2017rmpe},
which combines a person-detection front-end with a symmetric spatial transformer
and a parametric pose non-maximum suppression step.
Top-down methods generally have strong per-person keypoint localisation accuracy
due to the simplification of the pose estimation task to a single person at a time.
However, there are two main limitations to this approach.
Firstly, runtime scales linearly with the number of detected people.
Secondly, the performance of these methods is limited by the quality of the person detection stage:
any incorrectly detected or missed bounding boxes will lead to missed or incorrectly grouped keypoints,
with no opportunity for the pose estimator to recover from these errors.

\subsection{Bottom-up approaches}
\label{sec_lit_hpe_bottomup}

Bottom-up methods work in reverse.
The first task is to detect all the keypoints in an image
and then a separate grouping mechanism assigns each keypoint to a person.
Detecting all keypoints in a single pass allows for a runtime that is approximately constant in the number of people,
and detection errors are localised to individual keypoints rather than to entire person instances.
Grouping, being a critical step (specifically for this work), is reviewed further in Section~\ref{sec_lit_grouping}.
The exception to this is end-to-end transformer methods (Section~\ref{sec_lit_hpe_transformer}),
which do not require a separate grouping stage and are discussed separately.

\subsection{Transformer-based approaches}
\label{sec_lit_hpe_transformer}

The transformer architecture has been applied to HPE in two distinct ways:
as a drop-in replacement for the convolutional backbone of a standard heatmap-based estimator,
and as a complete reimagining of the multi-person task
using end-to-end set prediction with no separate grouping stage.
The latter approach cannot be decoupled from the backbone
and does not produce embeddings that a downstream grouping module could use.

TokenPose~\cite{li2021tokenpose} demonstrated that a pure-transformer
network where each keypoint is a learned token could match the
heatmap accuracy of CNN baselines at comparable parameter counts.
HRFormer~\cite{yuan2021hrformer} preserved the parallel-resolution
structure of HRNet while replacing convolutions with local-window
self-attention, recovering high-resolution feature representations.
ViTPose~\cite{xu2022vitpose} took the opposite design choice -
a plain, monolithic Vision Transformer as the backbone with no
high-resolution path and no domain-specific inductive bias for
keypoint localisation -
it then showed that this minimal architecture can scale cleanly
to model sizes from 100M up to 1B parameters (improving with scale).
A simple deconvolutional head turns the patch-token grid into
per-keypoint heatmaps, so ViTPose plugs into the same downstream pipeline as a CNN backbone.
Its successor ViTPose++~\cite{xu2023vitposepp} extends the original ViTPose by adding a
mixture-of-experts (MoE) formulation to cover the wholebody, hand, and
animal-pose distributions within a single trunk.

A separate line of transformer work reimagines multi-person HPE as direct set prediction,
so the keypoints are generated already grouped by person,
meaning there is no need for a separate grouping stage.
PETR~\cite{shi2022petr} was the first end-to-end pose detector to
apply the DETR~\cite{carion2020detr} set-prediction paradigm to human keypoints:
a fixed pool of pose queries cross-attends to image features,
and each query learns to predict a complete per-person skeleton (one shot).
The predictions form an unordered set with no fixed pairing to the ground-truth people,
and Hungarian matching at training time assigns each ground-truth person to one of the queries.
Here the binding between keypoints and person is the query identity.
QueryPose~\cite{xiao2022querypose} sparsified the attention pattern by
splitting each pose query into part-level sub-queries.
ED-Pose~\cite{yang2023edpose} and GroupPose~\cite{liu2023grouppose}
modified the query-based formulation with ED-Pose adding a person-box
and GroupPose simplifying the architecture to allow keypoints of the same person to communicate;
both reaching competitive performance on the COCO benchmark without any
post-hoc grouping step.

\subsection{Heatmap and regression trends}
\label{sec_lit_hpe_regression}

A long-standing issue in regression-based pose estimation is the non-differentiability of the argmax operator
which maps continuous outputs to discrete keypoint coordinates.
Sun et al.~\cite{sun2018integral} addressed this with integral pose regression,
introducing a differentiable soft-argmax that bridged the heatmap and regression paradigms.
The differentiable spatial-to-numerical transform (DSNT)~\cite{nibali2018dsnt}
later formalised the same idea as a normalised spatial-coordinate layer.

RLE~\cite{li2021rle} offered a more direct probabilistic treatment,
obtained by modelling residuals between predicted and ground-truth coordinates with a flow-based density,
and it was the first regression method to exceed heatmap accuracy on COCO.
SimCC~\cite{li2022simcc} took a different route,
reframing coordinate prediction as two-axis classification,
which became the basis for the real-time RTMPose family~\cite{jiang2023rtmpose}.
RTMO~\cite{lu2024rtmo} later unified detection and regression in a single stage,
and RTMW~\cite{jiang2024rtmw} extended this to whole-body pose at real-time rates.
The YOLO family follows the same single-stage pattern,
attaching a per-instance keypoint regression head to its object detector,
with YOLO11-pose~\cite{ultralytics2024yolo11} as the current representative.

For this context, the specific paradigm is not critical.
The grouping module proposed in Chapter~\ref{chap_3} operates on per-keypoint coordinates
and is indifferent to whether those coordinates came from a heatmap or a regressor.

\subsection{Foundation models for pose estimation}
\label{sec_lit_hpe_foundation}

Recent work has also moved human pose backbones towards foundation-model scale.
Sapiens~\cite{khirodkar2024sapiens}, for example,
pretrains Vision Transformer backbones ranging from 0.3 to 2 billion parameters
using 300M curated human images.
This produces a unified human-vision backbone that achieves state-of-the-art accuracy
across pose estimation, segmentation, depth, and surface normal prediction.

A similar move towards promptable and dataset-unified methods can be seen in X-Pose~\cite{yang2024xpose},
which detects arbitrary keypoint categories based on either a textual or visual prompt
across thirteen datasets.
PCT~\cite{geng2023pct} takes a different approach
using discrete codebook tokens to represent poses and improve robustness to occlusion.

These foundation-model approaches are relevant here mainly as context rather than direct points of comparison.
Larger pretrained backbones can reduce the localisation error of individual keypoints to near saturation.
This further makes the case for studying grouping as a separate, modular stage after detection.

\subsection{Benchmarks}
\label{sec_lit_hpe_benchmarks}

Three datasets are particularly common in the 2D multi-person pose literature.
MPII Human Pose~\cite{andriluka2014mpii} established much of the modern benchmark protocol,
with 25K images and approximately 40K annotated people across 410 activities.
However, its focus is mainly on single-person crops,
and it has largely been superseded for multi-person research.

COCO~\cite{lin2014coco} is now the de facto benchmark for multi-person pose estimation,
using seventeen keypoints per person
and containing approximately 250K annotated instances across the train, val, and test partitions.
The val2017 set is commonly used for reporting results
and is also the set used for the COCO evaluations in Chapter~\ref{chap_4}.

CrowdPose~\cite{li2019crowdpose}, on the other hand, is designed specifically for crowded scenes.
It includes a crowd-index metric and approximately 20K images selected for significant inter-person occlusion,
making it a useful benchmark for testing grouping methods under more difficult conditions,
although it is not used in the experiments later in this work.
The OCHuman benchmark, introduced alongside Pose2Seg~\cite{zhang2019pose2seg},
also focuses on heavily occluded people and is therefore relevant when considering grouping robustness.

Across these benchmarks,
average precision (AP) under the Object Keypoint Similarity (OKS) metric is the dominant evaluation measure.
OKS was introduced with the COCO benchmark~\cite{lin2014coco}
and measures the agreement between predicted and ground-truth keypoints using a Gaussian-based tolerance,
where the variance depends on the keypoint type and the size of the person's bounding box.
A prediction is counted as correct above an OKS threshold,
with AP normally averaged over thresholds in the range $[0.50, 0.95]$ in steps of $0.05$.

The earlier MPII-era metric, Percentage of Correct Keypoints (PCK), uses a different approach:
a keypoint is counted as correct when it falls within a fixed fraction of the head segment length.
PCK is still sometimes reported for MPII,
but AP@OKS has largely replaced it for multi-person benchmarks.

AP@OKS is not reported directly in Chapter~\ref{chap_4}.
Instead, the grouping contribution is evaluated using Hungarian-matched per-joint accuracy
(see Section~\ref{sec_meth_evaluation}),
which separates the grouping decision from the quality of the initial keypoint detections.
The grouping baselines discussed in this chapter are, however,
reported using AP@OKS in their original publications,
and the relationship between the two metrics is discussed further in Chapter~\ref{chap_4}.

\subsection{Pose tracking}
\label{sec_lit_hpe_tracking}

Multi-person 2D pose tracking extends per-image HPE by associating people across multiple frames.
PoseFlow~\cite{xiu2018poseflow} is a well-known online tracker
that maintains person identities using pose-similarity matching over a small temporal window,
while LightTrack~\cite{ning2020lighttrack} uses a top-down tracking framework with a learned re-identification head.
The PoseTrack benchmark~\cite{andriluka2018posetrack} is the standard evaluation set
for this video-based pose tracking problem.

Pose tracking is outside the scope of this work,
which is limited to grouping keypoints within individual images.
However, grouping remains an important step for each frame.
Once the per-frame skeletons have been assembled,
temporal association can generally be treated as a re-identification problem over a short sequence of frames.
The per-image grouping module developed in Chapter~\ref{chap_3} can therefore be used as a separate stage
before temporal association,
with standard ReID methods handling the subsequent tracking step.

\subsection{Position of this work within the HPE landscape}
\label{sec_lit_hpe_position}

The HPE literature has converged on a relatively small set of strong single-image detection backbones,
with per-keypoint localisation accuracy now reaching saturation.
HRNet-family backbones already achieve mean OKS-based average precision above $0.75$ on COCO val2017,
while more recent foundation-model approaches such as Sapiens~\cite{khirodkar2024sapiens}
improve on this further through pretraining on hundreds of millions of human-centric images.
For systems where models of this scale are not practical,
real-time approaches such as RTMPose~\cite{jiang2023rtmpose} and RTMO~\cite{lu2024rtmo}
achieve much of the same accuracy using considerably smaller models.
This suggests that, for many practical multi-person systems,
localising individual keypoints is no longer the main limitation on performance.

For bottom-up approaches, an important source of error therefore lies after detection,
particularly in deciding which detected keypoints belong to the same person.
Most of the established methods reviewed above
make this grouping decision in a way that is closely tied to the detector.
Part-affinity fields, for example, are predicted by the same network that generates the keypoint heatmaps.
Associative-embedding tags are similarly trained together with the heatmap loss on a particular backbone.
Centre-regression methods such as DEKR also combine detection and grouping
through a backbone-specific keypoint-from-centre offset head.
End-to-end transformer pose detectors couple the two stages even more closely:
in these models there is no separate per-keypoint representation to group,
since the keypoints are already associated through the pose query that produced them.
As a result, replacing the detector after training is impossible,
and grouping methods are often tied to the detector and grouping head used by the original architecture.

In comparison, the literature on \emph{modular} grouping is sparse.
Here, modular grouping refers to methods that accept detected keypoints from any upstream detector,
perform grouping as a separate stage,
and return a person identity for each keypoint.
Several recent methods can be viewed in this way.
HGG~\cite{jin2020hgg} treats the problem as higher-order graph matching between detected keypoints,
while GEC~\cite{maher2025gec} formulates it as binary edge classification.
The instance-aware method of Yang et al.~\cite{yang2023instanceaware}
instead uses a learned per-person assignment head.

The method proposed in this work follows the same modular direction.
PG-GAT, presented in Chapter~\ref{chap_3},
uses a graph neural network to operate on keypoints produced by any upstream detector
and generates an embedding for each detected keypoint.
These embeddings can then be clustered to recover the individual person identities.
The following sections review the literature that most directly supports this approach:
grouping mechanisms in Section~\ref{sec_lit_grouping}
and graph neural networks for pose in Section~\ref{sec_lit_gnn_pose}.